\documentclass[11pt, a4paper]{article}
\usepackage[T1]{fontenc}
\usepackage[utf8]{inputenc}
\usepackage[a4paper, top=2.5cm, bottom=2.5cm, left=2.5cm, right=2.5cm]{geometry}
\usepackage{booktabs}
\usepackage{tabularx}
\usepackage{enumitem}
\usepackage[hidelinks]{hyperref}
\usepackage{parskip}
\usepackage{graphicx}
\usepackage{amsmath}
\usepackage{float}
\setlength{\parindent}{0pt}
\pagestyle{plain}

\begin{document}

% ------------------------------------------------------------------------------
% PORTADA
% ------------------------------------------------------------------------------
\begin{center}
{\Large \textbf{Pontificia Universidad Javeriana}}\\[0.3cm]
Facultad de Ciencias Economicas y Administrativas\\[0.3cm]
Estadistica -- 2026-II\\[1cm]
{\LARGE \textbf{Meritocracia y origen territorial: ?`determina el lugar de origen el exito en educacion superior?}}\\[0.5cm]
{\large Informe Final del Proyecto}\\[1cm]
\textbf{Grupo 1}\\[0.3cm]
Perdigon Salinas Juan Felipe\\
Paez Sabogal Laura Valentina\\
Rodriguez Duque Catalina\\
Ramirez Cuesta Wendy Julieth\\[1cm]
27 de noviembre de 2026
\end{center}

\newpage
\tableofcontents
\newpage

% ------------------------------------------------------------------------------
\section{Introduccion}
% ------------------------------------------------------------------------------
% Contexto: desigualdad territorial en Colombia y acceso a educacion de calidad.
% Pregunta: un estudiante de Cauca, Choco o Narino con alto puntaje SB11
% obtiene el mismo resultado en SBPro que uno de Bogota?
% Concepto de meritocracia vs determinismo geografico.

% ------------------------------------------------------------------------------
\section{Datos y Metodologia}
% ------------------------------------------------------------------------------
% Fuente: cruce Saber 11 -- Saber Pro (ICFES).
% Archivos: cruce\_saber11\_saberpro.parquet, saber11.parquet, saber\_pro.parquet.
% Variable dependiente: PUNT\_GLOBAL\_sbpro.
% Predictores: PUNT\_GLOBAL\_sb11, COLE\_NATURALEZA, COLE\_AREA\_UBICACION,
%   region (Cauca/Choco/Narino vs Bogota), FAMI\_ESTRATOVIVIENDA\_sb11.
% Construccion de la variable region a partir de ESTU\_DEPTO\_RESIDE\_sb11.
% Modelo OLS completo y modelo sin SB11 (solo variables de origen).

% ------------------------------------------------------------------------------
\section{Analisis Exploratorio (Fase 1)}
% ------------------------------------------------------------------------------
% Distribucion de PUNT\_GLOBAL\_sbpro y PUNT\_GLOBAL\_sb11 por region.
% Tablas de frecuencias para COLE\_NATURALEZA y COLE\_AREA\_UBICACION.
% Estadisticas descriptivas por region y estrato.

% ------------------------------------------------------------------------------
\section{Intervalos de Confianza y Pruebas de Hipotesis (Fase 2)}
% ------------------------------------------------------------------------------
% IC para la media de PUNT\_GLOBAL\_sbpro por region.
% Prueba t para diferencia de medias Bogota vs Periferia\_sur.
% Prueba chi-cuadrado para asociacion entre region y COLE\_NATURALEZA.

% ------------------------------------------------------------------------------
\section{Modelo de Regresion (Fase 3)}
% ------------------------------------------------------------------------------
% Modelo completo: PUNT\_GLOBAL\_sbpro $\sim$ PUNT\_GLOBAL\_sb11 + COLE\_NATURALEZA +
%   COLE\_AREA\_UBICACION + region + FAMI\_ESTRATOVIVIENDA\_sb11.
% Modelo sin SB11: solo variables de origen territorial y socioeconomicas.
% Tabla de coeficientes con IC al 95\%.
% Comparacion de modelos: R2 ajustado, AIC, BIC.
% Predicciones: P90-SB11 rural Cauca vs P10-SB11 privado Bogota.

% ------------------------------------------------------------------------------
\section{Diagnosticos y Robustez}
% ------------------------------------------------------------------------------
% Graficos de residuales (residuales vs ajustados, QQ-plot, Scale-Location).
% Test de Breusch-Pagan (homocedasticidad).
% VIF (multicolinealidad).
% Distancia de Cook (observaciones influyentes).
% Errores robustos HC1: comparacion con errores estandar clasicos.

% ------------------------------------------------------------------------------
\section{Resultados y Discusion}
% ------------------------------------------------------------------------------
% Magnitud del efecto de region controlando por SB11.
% Cuanto explica el origen territorial vs el merito individual?
% Grafico de coeficientes, observado vs predicho, comparacion de predicciones.
% Reflexion sobre meritocracia y barreras estructurales.

% ------------------------------------------------------------------------------
\section{Conclusiones}
% ------------------------------------------------------------------------------
% Respuesta a la pregunta de investigacion.
% Limitaciones del estudio.
% Implicaciones para politica publica de equidad territorial en educacion.

% ------------------------------------------------------------------------------
\section*{Referencias}
% ------------------------------------------------------------------------------
% ICFES, literatura sobre desigualdad territorial y educacion en Colombia.

% ------------------------------------------------------------------------------
\section*{Anexo: Codigo R}
% ------------------------------------------------------------------------------
% Incluir codigo R utilizado (fase3\_grupo1.R) o referenciarlo.

\end{document}
